
\SetRegion{目录}{teal}{compass}
\phantomsection\label{ch1-contents}\pdfbookmark[0]{目录}{ch1-bm-contents}
\kicker{CONTENTS / 全书目录}
\PageTitle{目录}{参与、供给、复用、治理与公共价值}
\TOCLine{}{执行摘要}{\pageref{ch1-executive}}
\TOCLine{}{引言：开源共创智能未来}{\pageref{ch1-preface}}
\TOCLine{}{全书结构与阅读路径}{\pageref{ch1-architecture}}
\TOCLine{}{宣言结构与全文}{\pageref{ch1-declarationmap}}
\TOCHead{01}{从开放协作到公共价值}{teal}{发展方向与观察方法}
\TOCLine{1.1}{开源：共同创新与公共价值}{\pageref{ch1-ch1sec1}}
\TOCLine{1.2}{OSDGs：面向智能时代的开源发展目标}{\pageref{ch1-ch1sec2}}
\TOCLine{1.3}{OSDI：指标、证据与观察方法}{\pageref{ch1-ch1sec3}}
\TOCLine{}{第一章回望}{\pageref{ch1-closing}}
\TOCHead{02}{全球开源生态观测地图}{blue}{GOSO / 参与和资源分布在哪里}
\TOCLine{2.1}{全球开源活跃地图}{32}
\TOCLine{2.2}{全球开源贡献地图}{38}
\TOCLine{2.3}{全球开源项目影响力地图}{44}
\TOCLine{2.4}{全球开源人才地图}{50}
\TOCLine{2.5}{全球开源资产地图}{55}
\TOCLine{2.6}{全球开源多样性地图}{63}
\TOCLine{2.7}{全球开源治理与社会公益地图}{70}
\TOCHead{03}{全球开源影响力排行榜}{violet}{GOSR / 哪些主体形成何种影响}
\TOCLine{3.1}{全球行政区划影响力排行榜}{78}
\TOCLine{3.2}{全球开源企业影响力排行榜}{85}
\TOCLine{3.3}{全球开源基金会影响力排行榜}{92}
\TOCLine{3.4}{全球开源项目影响力排行榜}{97}
\TOCLine{3.5}{全球新势力开源仓库与项目排行榜}{108}
\TOCLine{3.6}{全球开源模型影响力排行榜}{112}
\clearpage
\kicker{CONTENTS / 全书目录}
\PageTitle{目录}{从技术结构，走向共同任务}
\TOCHead{04}{全球开源技术全景洞察}{rust}{GOSL / 技术如何组织与演进}
\TOCLine{4.1}{Agentic AI 全景洞察}{118}
\TOCLine{4.2}{AI 大模型全景洞察}{125}
\TOCLine{4.3}{云基础设施全景洞察}{132}
\TOCLine{4.4}{数据库系统全景洞察}{133}
\TOCLine{4.5}{RISC-V 全景洞察}{134}
\TOCLine{4.6}{PyTorch生态观察样本}{135}
\TOCHead{05}{全球开源大模型发展洞察}{green}{GOSM / 模型如何被供给、复用和开放}
\TOCLine{5.1}{供给与贡献}{139}
\TOCLine{5.2}{网络与影响力}{143}
\TOCLine{5.3}{开放条件与分级}{148}
\TOCHead{06}{未来展望与行动建议}{gold}{行动主体、支持条件与后续观察}
\TOCLine{6.1}{核心洞察与趋势展望}{155}
\TOCLine{6.2}{面向政府与国际组织的建议}{157}
\TOCLine{6.3}{面向企业与研究机构的建议}{159}
\TOCLine{6.4}{面向社区与全球伙伴的建议}{161}
\TOCLine{6.5}{结语与后续观察}{162}
\TOCHead{A}{附录与致谢}{teal}{共同概念、资料来源与协作贡献}
\TOCLine{A}{指标体系与数据来源}{164}
\TOCLine{}{致谢}{171}
\vspace{1mm}\noindent\includegraphics[height=9mm]{figures/new/dimension_strip.pdf}\par\src{D01—D10：十维导视，见第\pageref{ch1-dimensionoverview}页。页码与本PDF页面序号一致。A3竖向图页保留完整尺寸，可放大阅读。}


\clearpage

